% !TEX root = final_report.tex
\documentclass[final_report.tex]{subfiles}
\setcounter{chapter}{4}

\begin{document}
\ourchapter{The matched-budget pretext study}\label{ch:matched-budget}

\stub{Chapter preamble: with the signal established, this chapter measures what it
costs under controlled conditions, and produces the four converged encoders every
downstream measurement in \cref{ch:downstream} is run on.}

\section{Why matched budgets}\label{sec:why-matched-budgets}
\stub{The convergence confound: an arm that is merely slower looks like an arm that
is worse. The stand-down that followed, and what a matched budget has to rule out.}

\section{Design: four arms, one factor per step}\label{sec:four-arms}
\stub{Baseline, rotated pixels unsupervised, RoPART \(\pm 30^\circ\), RoPART
\(\pm 90^\circ\). Each step changes exactly one thing, which is what makes a later
downstream difference attributable to the \dphi{} target rather than to rotated
patches acting as augmentation.}

\section{Convergence and pretext performance}\label{sec:convergence}
\stub{Training behaviour across the four arms at matched budget, and confirmation
that each has converged rather than merely stopped.}

\section{The cost of rotation range}\label{sec:range-cost}
\stub{The intrinsic gates across \(\pm 30^\circ\) and \(\pm 90^\circ\). The gaps
survive budget matching, so the graded degradation is a property of the task ---
and, floor-normalised, the wider range is a trade rather than a failure.}

\section{The pretext-level no-degradation check}\label{sec:no-degradation-pretext}
\stub{Adding the \dphi{} target does not cost translation. States plainly that this
is a pretext-level result, on the task family the encoder was trained on, and so
does not settle the downstream question.}

\section{Discussion}\label{sec:matched-budget-discussion}
\stub{What the four arms can and cannot attribute, and which structural looseness
tracks rotated pixels rather than the target itself.}

\ifSubfilesClassLoaded{\pagebreak\printbibliography}{}
\end{document}
